\chapter{2011年广东卷（理）}

\section{单选题}

\begin{problem}
设复数 \(z\) 满足 \((1 + \i)z = 2\)，其中 \(\i\) 为虚数单位，则 \(z =\)（\quad）

\choices
    {\(1 + \i\)}
    {\(1 - \i\)}
    {\(2 + 2\i\)}
    {\(2 - 2\i\)}
\end{problem}

\begin{answer}
B
\end{answer}

\begin{solution}
两边同除以 \(1+\i\)，得 \(z=\dfrac{2}{1+\i}=\dfrac{2(1-\i)}{(1+\i)(1-\i)}=1-\i\)，故选 B.
\end{solution}

\begin{problem}
已知集合 \(A = \{(x, y) | x, y\) 为实数，且 \(x^2 + y^2 = 1\}\)，\(B = \{(x, y) | x, y\) 为实数，且 \(y = x\}\)，则 \(A \cap B\) 的元素个数为（\quad）

\choices
    {\(0\)}
    {\(1\)}
    {\(2\)}
    {\(3\)}
\end{problem}

\begin{answer}
C
\end{answer}

\begin{solution}
点 \((x,y)\) 同时属于 \(A\) 和 \(B\) 时，满足 \(y=x\) 以及 \(x^2+y^2=1\)，所以 \(2x^2=1\)，从而 \(x=\pm\frac{\sqrt{2}}{2}\)，且 \(y=x\).因此交点为 \(\left(\frac{\sqrt{2}}{2},\frac{\sqrt{2}}{2}\right)\) 和 \(\left(-\frac{\sqrt{2}}{2},-\frac{\sqrt{2}}{2}\right)\)，共有 \(2\) 个，故选 C.
\end{solution}

\begin{problem}
若非零向量 \(a\)，\(b\)，\(c\) 满足 \(a \parallel b\) 且 \(a \perp c\)，则 \(c \cdot (a + 2b) =\)（\quad）

\choices
    {\(4\)}
    {\(3\)}
    {\(2\)}
    {\(0\)}
\end{problem}

\begin{answer}
D
\end{answer}

\begin{solution}
因为 \(a\parallel b\)，且 \(a\perp c\)，所以 \(b\perp c\)，从而 \(c\cdot a=0\) 且 \(c\cdot b=0\).因此 \(c\cdot(a+2b)=c\cdot a+2c\cdot b=0\)，故选 D.
\end{solution}

\begin{problem}
设函数 \( f(x) \) 和 \( g(x) \) 分别是 \( \R \) 上的偶函数和奇函数，则下列结论恒成立的是（\quad）

\choices
    {\(f(x) + |g(x)|\) 是偶函数}
    {\(f(x) - |g(x)|\) 是奇函数}
    {\(|f(x)| + g(x)\) 是偶函数}
    {\( |f(x)| - g(x) \) 是奇函数}
\end{problem}

\begin{answer}
A
\end{answer}

\begin{solution}
因为 \(f(x)\) 是偶函数，\(g(x)\) 是奇函数，所以 \(f(-x)=f(x)\)，\(g(-x)=-g(x)\)，进而 \(\lvert g(-x)\rvert=\lvert g(x)\rvert\).于是 \(f(x)+\lvert g(x)\rvert\) 在 \(-x\) 处的值仍为 \(f(x)+\lvert g(x)\rvert\)，它恒为偶函数，故选 A.
\end{solution}

\begin{problem}
已知平面直角坐标系 \(xOy\) 上的区域 \(D\) 由不等式组 \(\begin{cases}
0 \leqslant x \leqslant \sqrt{2},\\
y \leqslant 2,\\
x \leqslant \sqrt{2} y
\end{cases}\) 给定. 若 \(M(x,y)\) 为 \(D\) 上的动点，点 \(A\) 的坐标为 \((\sqrt{2},1)\)，则 \(z = \overrightarrow{OM} \cdot \overrightarrow{OA}\) 的最大值为（\quad）

\choices
    {\(4\sqrt{2}\)}
    {\(3\sqrt{2}\)}
    {4}
    {3}
\end{problem}

\begin{answer}
C
\end{answer}

\begin{solution}
由区域条件知 \(0\leqslant x\leqslant\sqrt{2}\)，且 \(y\leqslant2\).又因为 \(z=\overrightarrow{OM}\cdot\overrightarrow{OA}=\sqrt{2}x+y\)，所以 \(z\leqslant\sqrt{2}\cdot\sqrt{2}+2=4\).点 \((\sqrt{2},2)\) 满足 \(x\leqslant\sqrt{2}y\)，属于 \(D\)，并且在该点取得上界，故最大值为 \(4\)，选 C.
\end{solution}

\begin{problem}
甲，乙两队进行排球决赛. 现在的情形是甲队只要再赢一局就获冠军，乙队需要再赢两局才能得冠军. 若两队胜每局的概率相同，则甲队获得冠军的概率为（\quad）

\choices
    {\(\frac{1}{2}\)}
    {\( \frac{3}{5} \)}
    {\( \frac{2}{3} \)}
    {\( \frac{3}{4} \)}
\end{problem}

\begin{answer}
D
\end{answer}

\begin{solution}
甲队在第一局获胜的概率为 \(\frac{1}{2}\).若第一局乙队获胜，甲队必须在第二局获胜，才能先达到所需的胜局数，这一事件的概率为 \(\frac{1}{2}\cdot\frac{1}{2}=\frac{1}{4}\).因此甲队获得冠军的概率为 \(\frac{1}{2}+\frac{1}{4}=\frac{3}{4}\)，故选 D.
\end{solution}

\begin{problem}
如图，某几何体的正视图（主视图）是平行四边形，侧视图（左视图）和俯视图都是矩形，则该几何体的体积为（\quad）

\bitmapfigure[max width=0.72\linewidth]{img/2011/guangdong_science/q7_fig1.png}

\choices
    {\(6\sqrt{3}\)}
    {\(9\sqrt{3}\)}
    {\(12\sqrt{3}\)}
    {\(18\sqrt{3}\)}
\end{problem}

\begin{answer}
B
\end{answer}

\begin{solution}
由三视图可知，该几何体可看作以正视图中的平行四边形为底面的斜四棱柱.俯视图表明平行四边形的上、下边在水平方向错开 \(1\)，而其斜边长为 \(2\)，所以平行四边形的高为 \(\sqrt{2^2-1^2}=\sqrt{3}\)，底面积为 \(3\sqrt{3}\).侧视图给出两底面之间的距离为 \(3\)，故几何体体积为 \(3\sqrt{3}\times3=9\sqrt{3}\)，选 B.
\end{solution}

\begin{problem}
设 \(S\) 是整数集 \(\Z\) 的非空子集，如果对于 \(\forall a,b\in S\)，有 \(ab \in S\)，则称 \(S\) 关于数的乘法是封闭的. 若 \(T\)，\(V\) 是 \(\Z\) 的两个不相交的非空子集，\(T \cup V = \Z\). 且对于 \(\forall a,b,c\in T\)，有 \(abc \in T\)，\(\forall x,y,z\in V\)，有 \(xyz \in V\). 则下列结论恒成立的是（\quad）

\choices
    {\(T\)，\( V\) 中至少有一个关于乘法是封闭的}
    {\(T\)，\(V\) 中至多有一个关于乘法是封闭的}
    {\(T\)，\(V\) 中有且只有一个关于乘法是封闭的}
    {\(T\)，\(V\) 中每一个关于乘法都是封闭的}
\end{problem}

\begin{answer}
A
\end{answer}

\begin{solution}
整数 \(1\) 必属于 \(T\) 或 \(V\) 中的一个.若 \(1\in T\)，则任取 \(a,b\in T\)，由题设关于三个元素的乘积的条件，得 \(ab=ab\cdot1\in T\)，所以 \(T\) 关于乘法封闭.若 \(1\in V\)，同理可得 \(V\) 关于乘法封闭.因此 \(T\)，\(V\) 中至少有一个关于乘法是封闭的，故选 A.
\end{solution}

\section{填空题}

\begin{problem}
不等式 \(|x + 1| - |x - 3| \geqslant 0\) 的解集是\fillinblank{}.
\end{problem}

\begin{answer}
\(\left[1,+\infty\right)\)
\end{answer}

\begin{solution}
两边都是非负数，所以原不等式等价于 \(\lvert x+1\rvert\geqslant\lvert x-3\rvert\).两边平方得 \((x+1)^2\geqslant(x-3)^2\)，即 \(8x\geqslant8\)，所以 \(x\geqslant1\).故解集为 \(\left[1,+\infty\right)\).
\end{solution}

\begin{problem}
\(x\left(x - \frac{2}{x}\right)^7\) 的展开式中，\(x^4\) 的系数是\fillinblank{}. （用数字作答）
\end{problem}

\begin{answer}
84
\end{answer}

\begin{solution}
由二项式定理，\(\left(x-\frac{2}{x}\right)^7\) 的通项为 \(\mathrm C_7^k x^{7-k}\left(-\frac{2}{x}\right)^k=\mathrm C_7^k(-2)^k x^{7-2k}\).乘以外面的 \(x\) 后，通项的幂次为 \(8-2k\).令 \(8-2k=4\)，得 \(k=2\)，所以 \(x^4\) 的系数为 \(\mathrm C_7^2(-2)^2=21\times4=84\).
\end{solution}

\begin{problem}
等差数列 \(\{a_{n}\}\) 前9项的和等于前4项的和. 若 \(a_{1} = 1\)，\(a_{k} + a_{4} = 0\)，则 \(k = \)\fillinblank{}.
\end{problem}

\begin{answer}
10
\end{answer}

\begin{solution}
设公差为 \(d\).由 \(S_9=S_4\) 及等差数列前 \(n\) 项和公式，得 \(\frac{9}{2}(2+8d)=\frac{4}{2}(2+3d)\)，从而 \(d=-\frac{1}{6}\).于是 \(a_4=1+3d=\frac{1}{2}\)，且 \(a_k=1+(k-1)d=\frac{7-k}{6}\).由 \(a_k+a_4=0\)，得 \(\frac{7-k}{6}+\frac{1}{2}=0\)，所以 \(k=10\).
\end{solution}

\begin{problem}
函数 \( f(x) = x^3 - 3x^2 + 1 \) 在 \( x =\fillinblank{} \) 处取得极小值.
\end{problem}

\begin{answer}
2
\end{answer}

\begin{solution}
\(f'(x)=3x^2-6x=3x(x-2)\).当 \(x<0\) 或 \(x>2\) 时，\(f'(x)>0\)；当 \(0<x<2\) 时，\(f'(x)<0\).因此函数在 \(x=2\) 处由减变增，取得极小值.
\end{solution}

\begin{problem}
某数学老师身高 \(176\mathrm{~cm}\)，他爷爷，父亲和儿子的身高分别是 \(173\mathrm{~cm}\)，\(170\mathrm{~cm}\) 和 \(182\mathrm{~cm}\). 因儿子的身高与父亲的身高有关，该老师用线性回归分析的方法预测他孙子的身高为\fillinblank{}\(\mathrm{cm}\).
\end{problem}

\begin{answer}
185
\end{answer}

\begin{solution}
把父亲身高作为自变量、儿子身高作为因变量，得到三组数据 \((173,170)\)，\((170,176)\)，\((176,182)\).其平均值分别为 \(\bar x=173\)，\(\bar y=176\).因此回归直线斜率为
\(\hat b=\dfrac{(173-173)(170-176)+(170-173)(176-176)+(176-173)(182-176)}{(173-173)^2+(170-173)^2+(176-173)^2}=1\)，截距为 \(\hat a=\bar y-\hat b\bar x=3\).回归方程为 \(\hat y=x+3\).老师的儿子身高为 \(182\mathrm{~cm}\)，所以预测孙子身高为 \(182+3=185\mathrm{~cm}\).
\end{solution}

\section{选考题：填空题}

\begin{problem}
已知两曲线参数方程分别为 \(\begin{cases}
x = \sqrt{5} \cos \theta,\\
y = \sin \theta,
\end{cases}\)（\(0 \leqslant \theta < \pi\)）和 \(\begin{cases}
x = \frac{5}{4} t^2,\\
y = t,
\end{cases}\)（\(t \in \R\)），它们的交点坐标为\fillinblank{}.
\end{problem}

\begin{answer}
\(\left(1,\frac{2\sqrt{5}}{5}\right)\)
\end{answer}

\begin{solution}
由第一条曲线的参数方程得 \(x^2=5\cos^2\theta=5(1-y^2)\)，所以其普通方程为 \(\frac{x^2}{5}+y^2=1\)，并且由 \(0\leqslant\theta<\pi\) 知 \(y\geqslant0\).第二条曲线满足 \(x=\frac54y^2\).代入第一条曲线的方程，得 \(\frac{1}{5}\left(\frac54y^2\right)^2+y^2=1\)，即 \(5y^4+16y^2-16=0\).令 \(u=y^2\)，则 \(5u^2+16u-16=0\)，解得 \(u=\frac45\) 或 \(u=-4\).因 \(y\geqslant0\)，所以 \(y=\frac{2\sqrt{5}}{5}\)，进而 \(x=\frac54y^2=1\).故交点坐标为 \(\left(1,\frac{2\sqrt{5}}{5}\right)\).
\end{solution}

\begin{problem}
如图，过圆 \(O\) 外一点 \(P\) 分别作圆的切线和割线，交于 \(A\)，\(B\)，且 \(PB = 7\)，\(C\) 是圆上一点，使得 \(BC = 5\)，\(\angle BAC = \angle APB\)，则 \(AB =\fillinblank{}\).

\FigureLayoutDeclare{img/2011/guangdong_science/q15_fig1.png}{6.5cm}{33bp 43bp 56bp 0bp}
\bitmapfigure[max width=0.36\linewidth]{img/2011/guangdong_science/q15_fig1.png}
\end{problem}

\begin{answer}
\(\sqrt{35}\)
\end{answer}

\begin{solution}
由切线与弦所成的角等于弦所对的圆周角，得 \(\angle PAB=\angle ACB\).结合题设 \(\angle APB=\angle BAC\)，可知 \(\triangle PAB\sim\triangle ACB\).对应边给出 \(\dfrac{PB}{AB}=\dfrac{AB}{BC}\)，所以 \(AB^2=PB\cdot BC=7\times5=35\).由于长度为正，故 \(AB=\sqrt{35}\).
\end{solution}

\section{解答题}

\begin{problem}
\begin{enumerate}
\item 已知函数 \(f(x) = 2\sin \left(\frac{1}{3} x - \frac{\pi}{6}\right)\)，\(x \in \R\). 求 \( f\left(\frac{5\pi}{4}\right) \) 的值.
\item 设 \(\alpha\)，\(\beta \in \left[0, \frac{\pi}{2}\right]\)，\(f\left(3\alpha + \frac{\pi}{2}\right) = \frac{10}{13}\)，\(f(3\beta + 2\pi) = \frac{6}{5}\)，求 \(\cos (\alpha + \beta)\) 的值.
\end{enumerate}
\end{problem}

\begin{solution}
\begin{enumerate}
\item \(f\left(\frac{5\pi}{4}\right)=2\sin\left(\frac{5\pi}{12}-\frac{\pi}{6}\right)=2\sin\frac{\pi}{4}=\sqrt{2}\).
\item 由 \(f\left(3\alpha+\frac{\pi}{2}\right)=2\sin\alpha=\frac{10}{13}\)，得 \(\sin\alpha=\frac{5}{13}\).因 \(\alpha\in\left[0,\frac{\pi}{2}\right]\)，所以 \(\cos\alpha=\frac{12}{13}\).又由 \(f(3\beta+2\pi)=2\sin\left(\beta+\frac{\pi}{2}\right)=2\cos\beta=\frac{6}{5}\)，得 \(\cos\beta=\frac35\)，从而 \(\sin\beta=\frac45\).因此 \(\cos(\alpha+\beta)=\cos\alpha\cos\beta-\sin\alpha\sin\beta=\frac{12}{13}\cdot\frac35-\frac5{13}\cdot\frac45=\frac{16}{65}\).
\end{enumerate}
\end{solution}

\begin{problem}
为了解甲，乙两厂的产品质量，采用分层抽样的方法从甲，乙两厂生产的产品中分别抽取 14 件和 5 件，测量产品中微量元素 \(x\)，\(y\) 的含量（单位：毫克）. 下表是乙厂的 5 件产品的测量数据：

\begin{center}
\begin{tabular}{|c|c|c|c|c|c|}
\hline
编号 & 1 & 2 & 3 & 4 & 5 \\
\hline
\(x\) & 169 & 178 & 166 & 175 & 180 \\
\hline
\(y\) & 75 & 80 & 77 & 70 & 81 \\
\hline
\end{tabular}
\end{center}

\begin{enumerate}
\item 已知甲厂生产的产品共有 98 件，求乙厂生产的产品数量.
\item 当产品中的微量元素 \(x\)，\(y\) 满足 \(x \geqslant 175\) 且 \(y \geqslant 75\) 时，该产品为优等品. 用上述样本数据估计乙厂生产的优等品的数量.
\item 从乙厂抽出的上述 5 件产品中，随机抽取 2 件，求抽取的 2 件产品中优等品数 \(\xi\) 的分布列及其均值（即数学期望）.
\end{enumerate}
\end{problem}

\begin{solution}
\begin{enumerate}
\item 分层抽样中各层的抽样比相同.设乙厂有 \(N\) 件产品，则 \(\frac{14}{98}=\frac{5}{N}\)，解得 \(N=35\).
\item 样本中同时满足 \(x\geqslant175\) 和 \(y\geqslant75\) 的是第 2 件和第 5 件，共 2 件，样本频率为 \(\frac25\).因此估计乙厂优等品数量为 \(35\times\frac25=14\) 件.
\item 5 件产品中有 2 件优等品和 3 件非优等品.随机抽取 2 件时，\(\xi\) 的可能取值为 \(0,1,2\)，且
\begin{center}
\begin{tabular}{c|ccc}
\hline
\(\xi\) & 0 & 1 & 2 \\
\hline
\(P\) & \(\dfrac{\mathrm C_3^2}{\mathrm C_5^2}=\dfrac3{10}\) & \(\dfrac{\mathrm C_2^1\mathrm C_3^1}{\mathrm C_5^2}=\dfrac35\) & \(\dfrac{\mathrm C_2^2}{\mathrm C_5^2}=\dfrac1{10}\) \\
\hline
\end{tabular}
\end{center}
所以 \(E(\xi)=0\cdot\frac3{10}+1\cdot\frac35+2\cdot\frac1{10}=\frac45\).
\end{enumerate}
\end{solution}

\begin{problem}
如图，在锥体 \(P - ABCD\) 中，\(ABCD\) 是边长为 1 的菱形，且 \(\angle DAB = 60^{\circ}\)，\(PA = PD = \sqrt{2}\)，\(PB = 2\)，\(E\)，\(F\) 分别是 \(BC\)，\(PC\) 的中点.

\begin{enumerate}
\item 证明：\(AD \perp\) 平面 \(DEF\).
\item 求二面角 \(P-AD-B\) 的余弦值.
\end{enumerate}

\bitmapfigure[max width=0.36\linewidth]{img/2011/guangdong_science/q18_fig1.png}
\end{problem}

\begin{solution}
建立空间直角坐标系，令 \(A=(0,0,0)\)，\(B=(1,0,0)\)，\(D=\left(\frac12,\frac{\sqrt3}{2},0\right)\)，则 \(C=\left(\frac32,\frac{\sqrt3}{2},0\right)\).设 \(P=(u,v,w)\).由 \(PA=PD\) 得 \(u+\sqrt3v=1\)，由 \(PA=PB\) 得 \(u=-\frac12\)，所以 \(v=\frac{\sqrt3}{2}\).再由 \(PA^2=2\) 得 \(w^2=1\).图形关于底面反射不影响结论，取 \(w=1\)，即 \(P=\left(-\frac12,\frac{\sqrt3}{2},1\right)\).于是 \(E=\left(\frac54,\frac{\sqrt3}{4},0\right)\)，\(F=\left(\frac12,\frac{\sqrt3}{2},\frac12\right)\).

\begin{enumerate}
\item 有 \(\overrightarrow{AD}=\left(\frac12,\frac{\sqrt3}{2},0\right)\)，\(\overrightarrow{DE}=\left(\frac34,-\frac{\sqrt3}{4},0\right)\)，\(\overrightarrow{DF}=\left(0,0,\frac12\right)\).因此 \(\overrightarrow{AD}\cdot\overrightarrow{DE}=\frac38-\frac38=0\)，且 \(\overrightarrow{AD}\cdot\overrightarrow{DF}=0\).由于 \(DE\) 与 \(DF\) 是平面 \(DEF\) 内相交的两条直线，所以 \(AD\perp\) 平面 \(DEF\).
\item 设 \(G\) 为 \(AD\) 的中点.因为 \(PA=PD\)，所以 \(PG\perp AD\)；又因为 \(\triangle ABD\) 是边长为 \(1\) 的等边三角形，所以 \(BG\perp AD\).故 \(\angle PGB\) 是二面角 \(P-AD-B\) 的平面角.由 \(AG=\frac12\)，得 \(PG^2=PA^2-AG^2=2-\frac14=\frac74\)，\(BG^2=AB^2-AG^2=1-\frac14=\frac34\).在 \(\triangle PGB\) 中，\(PB=2\)，由余弦定理得
\[
\cos\angle PGB=\frac{PG^2+BG^2-PB^2}{2PG\cdot BG}=\frac{\frac74+\frac34-4}{2\cdot\frac{\sqrt7}{2}\cdot\frac{\sqrt3}{2}}=-\frac{\sqrt{21}}{7}.
\]
\end{enumerate}
\end{solution}

\begin{problem}
\begin{enumerate}
\item 设圆 \(C\) 与两圆 \((x + \sqrt{5})^2 + y^2 = 4\)，\((x - \sqrt{5})^2 + y^2 = 4\) 中的一个内切，另一个外切. 求 \(C\) 的圆心轨迹 \(L\) 的方程.
\item 已知点 \(M\left(\frac{3\sqrt{5}}{5}, \frac{4\sqrt{5}}{5}\right)\)，\(F(\sqrt{5}, 0)\)，且 \(P\) 为 \(L\) 上动点，求 \(\|MP\| - \|FP\|\) 的最大值及此时点 \(P\) 的坐标.
\end{enumerate}
\end{problem}

\begin{solution}
\begin{enumerate}
\item 设圆 \(C\) 的圆心为 \(C_0\)，半径为 \(r\)，两定圆的圆心分别为 \(F_1=(-\sqrt5,0)\) 和 \(F_2=(\sqrt5,0)\)，半径均为 \(2\).内切关系给出两圆心距为 \(\lvert r-2\rvert\)，外切关系给出另一圆心距为 \(r+2\).若 \(0<r\leqslant2\)，则两圆心距之和为 \(4\)，但 \(F_1F_2=2\sqrt5>4\)，不可能.因此 \(r>2\)，从而 \(\bigl\lvert C_0F_1\bigr\rvert\) 与 \(\bigl\lvert C_0F_2\bigr\rvert\) 之差的绝对值为 \(4\).所以轨迹是以 \(F_1,F_2\) 为焦点、实轴长为 \(4\) 的双曲线.其半焦距为 \(c=\sqrt5\)，半实轴长为 \(a=2\)，故 \(b^2=c^2-a^2=1\)，轨迹方程为 \(\dfrac{x^2}{4}-y^2=1\).
\item 点 \(M\) 与 \(F\) 的距离为 \(\lVert MF\rVert=2\)，直线 \(MF\) 的方程为 \(y=-2(x-\sqrt5)\).将其与 \(L\) 联立，得 \(15x^2-32\sqrt5x+84=0\)，解得 \(x=\frac{6\sqrt5}{5}\) 或 \(x=\frac{14\sqrt5}{15}\).代入直线方程，得到两个交点分别为 \(P_1=\left(\frac{6\sqrt5}{5},-\frac{2\sqrt5}{5}\right)\) 和 \(P_2=\left(\frac{14\sqrt5}{15},\frac{2\sqrt5}{15}\right)\).沿直线的横坐标顺序为 \(M\)，\(P_2\)，\(F\)，\(P_1\)，所以 \(F\) 在线段 \(MP_1\) 上.对任意 \(P\in L\)，三角不等式给出 \(\lVert MP\rVert-\lVert FP\rVert\leqslant\lVert MF\rVert=2\).在 \(P=P_1\) 时取等号，故最大值为 \(2\)，此时 \(P=\left(\frac{6\sqrt5}{5},-\frac{2\sqrt5}{5}\right)\).
\end{enumerate}
\end{solution}

\begin{problem}
\begin{enumerate}
\item 设 \(b > 0\)，数列 \(\{a_{n}\}\) 满足 \(a_{1} = b\)，\(a_{n} = \frac{nba_{n-1}}{a_{n-1} + 2n - 2}\)（\(n \geqslant 2\)）. 求数列 \(\{a_{n}\}\) 的通项公式.
\item 证明：对于一切正整数 \(n\)，\(a_{n} \leqslant \frac{b^{n+1}}{2^{n+1}} + 1\).
\end{enumerate}
\end{problem}

\begin{solution}
令 \(u_n=\frac{a_n}{n}\)，则由递推式得 \(u_1=b\)，且 \(u_n=\frac{b u_{n-1}}{u_{n-1}+2}\)（\(n\geqslant2\)）.

当 \(b=2\) 时，归纳可得 \(u_n=\frac2n\)：当 \(n=1\) 时成立，若 \(u_{n-1}=\frac{2}{n-1}\)，则 \(u_n=\frac{2\cdot\frac2{n-1}}{\frac2{n-1}+2}=\frac2n\).所以此时 \(a_n=2\).

当 \(b\neq2\) 时，先由 \(u_1=b=\frac{b(2-b)}{2-b}\) 知公式对 \(n=1\) 成立.若 \(u_{n-1}=\frac{b^{n-1}(2-b)}{2^{n-1}-b^{n-1}}\)，则
\[
u_n=\frac{b\,u_{n-1}}{u_{n-1}+2}=\frac{b^n(2-b)}{b^{n-1}(2-b)+2(2^{n-1}-b^{n-1})}=\frac{b^n(2-b)}{2^n-b^n}.
\]
因此
\[
a_n=\begin{cases}
2,&b=2,\\
\dfrac{n b^n(2-b)}{2^n-b^n},&b\neq2.
\end{cases}
\]

再证不等式.令 \(t=\frac b2>0\).当 \(b\neq2\) 时，利用等比数列求和可将上式写成 \(a_n=\dfrac{2nt^n}{1+t+t^2+\cdots+t^{n-1}}\)；当 \(b=2\) 时该表达式也给出 \(a_n=2\).由算术平均数不小于几何平均数，得 \(1+t+\cdots+t^{n-1}\geqslant n t^{(n-1)/2}\)，所以 \(a_n\leqslant2t^{(n+1)/2}\).又由基本不等式 \(1+t^{n+1}\geqslant2t^{(n+1)/2}\)，故 \(a_n\leqslant1+t^{n+1}=1+\frac{b^{n+1}}{2^{n+1}}\)，结论成立.
\end{solution}

\begin{problem}
\begin{enumerate}
\item 在平面直角坐标系 \(xOy\) 上，给定抛物线 \(L: y = \frac{1}{4} x^2\). 实数 \(p\)，\(q\) 满足 \(p^2 - 4q \geqslant 0\)，\(x_1\)，\(x_2\) 是方程 \(x^2 - px + q = 0\) 的两根. 记 \(\varphi(p, q) = \max \{|x_1|, |x_2|\}\). 过点 \(A\left(p_0, \frac{1}{4} p_0^2\right)\)（\(p_0 \neq 0\)）作 \(L\) 的切线交 \(y\) 轴于点 \(B\). 证明：对线段 \(AB\) 上的任一点 \(Q(p, q)\)，有 \(\varphi(p, q) = \frac{|p_0|}{2}\).
\item 设 \(M(a, b)\) 是定点，其中 \(a\)，\(b\) 满足 \(a^2 - 4b > 0\)，\(a \neq 0\). 过点 \(M(a, b)\) 作 \(L\) 的两条切线 \(l_1\)，\(l_2\)，切点分别为 \(E\left(p_1, \frac{1}{4} p_1^2\right)\)，\(E'\left(p_2, \frac{1}{4} p_2^2\right)\)，\(l_1\)，\(l_2\) 与 \(y\) 分别交于 \(F\)，\(F'\). 线段 \(EF\) 上异于两端点的点集记为 \(X\). 证明：\(M(a, b) \in X \Leftrightarrow |p_1| > |p_2| \Leftrightarrow \varphi(a, b) = \frac{|p_1|}{2}\).
\item 设 \(D = \left\{(x, y) \mid y \leqslant x - 1, y \geqslant \frac{1}{4} (x + 1)^2 - \frac{5}{4}\right\}\). 当点 \((p, q)\) 取遍 D 时，求 \(\varphi(p, q)\) 的最小值（记为 \(\varphi_{\min}\)）和最大值（记为 \(\varphi_{\max}\)）.
\end{enumerate}
\end{problem}

\begin{solution}
\begin{enumerate}
\item 抛物线在 \(A\) 处的切线斜率为 \(\frac{p_0}{2}\)，所以切线方程为 \(y=\frac{p_0}{2}x-\frac{p_0^2}{4}\)，从而 \(B=\left(0,-\frac{p_0^2}{4}\right)\).任取线段 \(AB\) 上的点 \(Q(p,q)\)，有 \(q=\frac{p_0}{2}p-\frac{p_0^2}{4}\)，且 \(p\) 位于 \(0\) 与 \(p_0\) 之间.于是
\[
x^2-px+q=x^2-px+\frac{p_0}{2}p-\frac{p_0^2}{4}=\left(x-\frac{p_0}{2}\right)\left(x-p+\frac{p_0}{2}\right).
\]
方程的两个根为 \(\frac{p_0}{2}\) 和 \(p-\frac{p_0}{2}\).因为 \(p\) 位于 \(0\) 与 \(p_0\) 之间，所以 \(\left\lvert p-\frac{p_0}{2}\right\rvert\leqslant\frac{|p_0|}{2}\).因此 \(\varphi(p,q)=\frac{|p_0|}{2}\).
\item 抛物线在参数为 \(p_i\) 的切点处的切线方程为 \(y=\frac{p_i}{2}x-\frac{p_i^2}{4}\).因其经过 \(M(a,b)\)，故 \(p_1,p_2\) 是方程 \(p^2-2ap+4b=0\) 的两个根，即 \(p_1+p_2=2a\).直线 \(EF\) 的端点横坐标分别为 \(p_1\) 和 \(0\)，而点 \(M\) 的横坐标为 \(a=\frac{p_1+p_2}{2}\).

当 \(p_1>0\) 时，\(M\) 在线段 \(EF\) 内部等价于 \(0<a<p_1\)，即 \(-p_1<p_2<p_1\)，等价于 \(|p_2|<|p_1|\).当 \(p_1<0\) 时，\(M\) 在线段 \(EF\) 内部等价于 \(p_1<a<0\)，同样等价于 \(|p_2|<|p_1|\).若 \(p_1=0\)，则由 \(a\neq0\) 有 \(p_2=2a\neq0\)，三者均不成立.因此 \(M(a,b)\in X\Longleftrightarrow |p_1|>|p_2|\).

又由 \(p_1+p_2=2a\) 以及 \(p_1p_2=4b\)，方程 \(x^2-ax+b=0\) 的两个根正是 \(\frac{p_1}{2}\) 和 \(\frac{p_2}{2}\).所以 \(\varphi(a,b)=\frac12\max\{|p_1|,|p_2|\}\)，从而 \(|p_1|>|p_2|\Longleftrightarrow\varphi(a,b)=\frac{|p_1|}{2}\).
\item 区域 \(D\) 非空要求 \(\frac14(p+1)^2-\frac54\leqslant p-1\)，化简得 \(0\leqslant p\leqslant2\).又因 \(q\leqslant p-1\)，有 \(p^2-4q\geqslant p^2-4p+4=(p-2)^2\geqslant0\).在固定的 \(p\in[0,2]\) 下，\(\varphi(p,q)=\frac{p+\sqrt{p^2-4q}}{2}\)，并且它随 \(q\) 增大而减小.

当 \(q=p-1\) 时，\(\sqrt{p^2-4q}=2-p\)，所以 \(\varphi(p,q)=1\)，从而 \(\varphi_{\min}=1\).当 \(q=\frac14(p+1)^2-\frac54\) 时，令 \(t=\sqrt{4-2p}\)，则 \(0\leqslant t\leqslant2\)，且 \(p=2-\frac{t^2}{2}\).此时
\[
\varphi(p,q)=\frac{p+t}{2}=1+\frac{t}{2}-\frac{t^2}{4}=\frac54-\frac{(t-1)^2}{4}\leqslant\frac54.
\]
当 \(t=1\)，即 \(p=\frac32\)，\(q=\frac5{16}\) 时取等号.因此 \(\varphi_{\max}=\frac54\).
\end{enumerate}
\end{solution}
